% -*- coding: utf-8 -*-
\chapter{国内外研究现状与理论基础}
\label{chap:related_work}

\section{通感一体化系统三大演进范式}
通感一体化系统依据其核心驱动目标与波形合成机制的演进历程，在学术界与工业界主要形成以下三大设计范式：

\subsection{通信优先感知体系（Comm-Centric ISAC）}
通信优先设计以标准移动蜂窝系统（如 3GPP 5G-Advanced 与 6G 候选波形）为基准框架。在该体制下，系统物理层依然沿用正交频分复用（OFDM）或离散傅里叶变换扩展正交频分复用（DFT-s-OFDM）信号结构，主要通过复用现有的解调参考信号（DMRS）或定位参考信号（PRS）进行雷达回波参数估计。接收机通过对已知发射符号与反射回波执行二维相关运算或匹配滤波，提取目标的距离、角度与径向多普勒速度。通信优先架构的优势在于能够最大程度兼容现有商用蜂窝移动通信协议与射频芯片，硬件升级成本低；然而其局限在于通信信号的高峰均比特性恶化了发射机功率放大器的直流转换效率，且导频子载波在频域的梳状或离散排布会导致雷达模糊函数产生严重的栅瓣和距离模糊。

\subsection{雷达优先通信体系（Radar-Centric ISAC）}
雷达优先设计以成熟的雷达体制（如线性调频 LFM 或调频连续波 FMCW）为载体，通过在雷达脉冲重复周期、初始相位或调频斜率上调制通信数据。例如，利用脉冲位置调制（PPM）或准正交啁啾分集实现控制信令广播。该架构能够保持雷达高发射功率、大动态范围与理想自相关特性的优势，硬件射频器件成熟；然而其核心缺陷在于通信数据传输速率受限于脉冲重复频率（PRF），难以满足低空智联网中高清视频图传与海量传感器数据回传所需的吉比特每秒（Gbps）级吞吐量需求。

\subsection{双功能一体化深度融合体系（DFRC）}
双功能雷达通信一体化（Dual-Functional Radar-Communication, DFRC）是当前通感一体化领域的核心前沿方向。该体制不再拘泥于单一传统波形的表层修补，而是从多天线阵列信号处理与联合优化理论出发，设计统一的基带发射信号与多波束预编码矩阵。发射信号在空间域同时合成指向下行通信移动用户的定向通信波束，以及覆盖待检测空域的广域雷达探测波束。通过在多目标空间滤波层解耦通信互干扰与雷达回波泄露，DFRC 系统能够在信息传输速率与目标探测克拉美罗界（CRB）之间灵活调节权衡权重，达到理论物理边界的最佳平衡。

\section{低空空地信道特性与高动态多普勒效应}
低空无人机系统所处的无线电磁传播环境介于传统地面蜂窝信道与高空卫星信道之间，具有独特的信道统计特性。

\subsection{空地三维几何多径信道表征}
设无人机飞行于高度为 $H$ 的空域，地面分布有 $M$ 个协同基站。基站与无人机之间的空间距离可表示为欧氏几何距离 $d_u$。空地视距（Line-of-Sight, LoS）链路的出现概率取决于仰角 $\theta$ 与地面建筑密度分布参数：
\begin{equation}
  P_{\mathrm{LoS}}(\theta) = \frac{1}{1 + a \exp\left(-b[\theta - a]\right)}
\end{equation}
其中 $a$ 与 $b$ 为表征典型城区、郊区或开阔地貌的环境常量。相应的空地信道 $\mathbf{H}_c \in \mathbb{C}^{K \times N_t}$ 可采用莱斯小尺度衰落模型表征：
\begin{equation}
  \mathbf{H}_c = \sqrt{\frac{K_R}{K_R + 1}}\mathbf{H}_{c,\mathrm{LoS}} + \sqrt{\frac{1}{K_R + 1}}\mathbf{H}_{c,\mathrm{NLoS}}
\end{equation}
其中 $K_R$ 为莱斯因子，$\mathbf{H}_{c,\mathrm{LoS}}$ 与 $\mathbf{H}_{c,\mathrm{NLoS}}$ 分别代表确定性视距分量与复高斯散射分量。

\subsection{高动态时变多普勒频移扩展}
当无人机以飞行速度向量 $\mathbf{v}_u$ 机动飞行时，其相对于地面基站或探测目标的相对速度产生显著的多普勒频偏：
\begin{equation}
  f_{D} = \frac{f_c}{c} \|\mathbf{v}_u\| \cos\psi
\end{equation}
其中 $f_c$ 为载波中心频率，$c$ 为光速，$\psi$ 为速度向量与波束传播方向的夹角。在 5G 毫米波频段（例如 $f_c = 28$\,GHz），当无人机巡航速度达到 $V_{\max} = 25$\,m/s 时，最大单向多普勒频移超过 2.3\,kHz，双向雷达回波多普勒频偏则高达 4.6\,kHz。若子载波间隔 $\Delta f = 15$\,kHz 或 30\,kHz，多普勒频偏占子载波间隔的比例显著提升，破坏子载波正交积分条件，导致严重的能量泄漏与信干噪比下降。

\section{现有典型波形与波束赋形技术对比分析}
为了明确现有各类技术方案的性能边界，表~\ref{tab:related_work_comparison} 给出了当前具有代表性的通感一体化设计方案的横向技术特征对比。

\begin{table}[htbp]
  \centering
  \caption{国内外现有典型通感一体化波形与资源管理方案特征对比}
  \label{tab:related_work_comparison}
  \small
  \resizebox{\textwidth}{!}{%
  \begin{tabular}{lllll}
    \toprule
    \textbf{文献方案} & \textbf{波形与体制架构} & \textbf{感知评估准则} & \textbf{通信性能指标} & \textbf{优化求解机制} \\
    \midrule
    顾等人方案\cite{gu2024papr} & 低 PAPR 多载波波形 & 雷达模糊度与峰值副瓣比 & 误码率与信干噪比 & 凸松弛与交替投影 \\
    万等人方案\cite{wan2024orthogonal} & 正交线性调频多分 (OCDM) & 距离-多普勒二维分辨率 & 接收信噪比与块错误率 & 菲涅尔变换匹配滤波 \\
    霍等人方案\cite{huo2024waveform} & MIMO-OFDM 探测波形 & 奈曼-皮尔逊检测概率 & 下行链路香农和速率 & 半定松弛与流形规划 \\
    仲等人方案\cite{zhong2022constant} & 恒模单载波雷达波形 & 脉冲压缩自相关旁瓣电平 & 多用户共信道干扰容限 & 黎曼流形共轭梯度法 \\
    付等人方案\cite{fu2024joint} & 无人机辅助联合预编码 & 目标测角克拉美罗下界 (CRB) & 加权系统和速率 & 逐次凸逼近与块坐标下降 \\
    左等人方案\cite{zuo2025robust} & 统计鲁棒发射波束赋形 & 最差工况感知回波信噪比 & 保密中断概率与保密速率 & S-程序变换与半正定松弛 \\
    \bottomrule
  \end{tabular}%
  }
\end{table}

\section{本章小结}
本章系统梳理了通感一体化系统的三大演进范式，深入剖析了低空空地信道的几何遮挡、莱斯衰落与毫米波高动态多普勒频移瓶颈，并对国内外现有主流波形体制与优化方案进行了多维对比。上述分析表明，在低空复杂多径与强机动环境下，设计一种兼顾低 PAPR、抗多普勒扩展并在数学上具备严谨 Pareto 最优性保证的双功能多载波波形，是低空智联网亟需攻克的理论瓶颈。
